\subsection{PRIMITIVE ELEMENTS OF A COGEBRA}

Throughout this no. we consider a cogebra E (Algebra, Chapter III, § 11, no. 1) with coproduct

\[
c \colon \mathrm{E} \rightarrow \mathrm{E} \otimes \mathrm{E}
\]

and counit \(\varepsilon\) (loc. cit., no. 2). Recall that \(\varepsilon\) is a linear form on the K-module E such that (with the canonical identification of \(\mathrm{E} \otimes \mathrm{K}\) and \(\mathrm{K} \otimes \mathrm{E}\) with E):

\[
\mathrm{Id} _ {\mathrm{E}} = (\varepsilon \otimes \mathrm{Id} _ {\mathrm{E}}) \circ c = (\mathrm{Id} _ {\mathrm{E}} \otimes \varepsilon) \circ c.
\]

Let \(E^{+}\) denote the kernel of \(\varepsilon\) and let u be an element of E such that

\[
c (u) = u \otimes u \quad \text { and } \quad \varepsilon (u) = 1.
\]

The K-module E is the direct sum of \(E^{+}\) and the submodule K.u which is free with basis u; let \(\pi_{u}: E \to E^{+}\) and \(\eta_{u}: E \to K.u\) denote the projectors associated with this decomposition. Then

\[
\pi_ {u} (x) = x - \varepsilon (x). u, \quad \eta_ {u} (x) = \varepsilon (x). u.\tag{1}
\]

DEFINITION 1. An element \(x\) of \(\mathbf{E}\) is called \(u\) -primitive if

\[
c (x) = x \otimes u + u \otimes x.\tag{2}
\]

The \(u\) -primitive elements of \(\mathbf{E}\) form a submodule of \(\mathbf{E}\) , denoted by \(\mathrm{P}_{u}(\mathrm{E})\) .

PROPOSITION 1. Every \(u\) -primitive element of \(\mathbf{E}\) belongs to \(\mathbf{E}^{+}\) .

\begin{enumerate}
\def\labelenumi{(\arabic{enumi})}
\setcounter{enumi}{1}
\tightlist
\item
  implies \(x = \varepsilon(x).u + \varepsilon(u).x = \varepsilon(x).u + x\) , whence \(\varepsilon(x) = 0\) .
\end{enumerate}

Remark. If \(x \in \mathrm{E}\) and \(c(x) = x' \otimes u + u \otimes x''\) , where \(x', x''\) are in \(\mathrm{E}^+\) , then \(x = \varepsilon(x')\) . \(u + \varepsilon(u)\) . \(x'' = x''\) ; similarly \(x = x'\) and \(x\) is \(u\) -primitive.

For all \(x \in \mathrm{E}^{+}\) , we write

\[
c _ {u} ^ {+} (x) = c (x) - x \otimes u - u \otimes x.\tag{3}
\]

PROPOSITION 2. We have

\[
\left(\pi_ {u} \otimes \pi_ {u}\right) \circ c = c _ {u} ^ {+} \circ \pi_ {u}.\tag{4}
\]

Let \(x\) be in E; then

\[
\begin{array}{r l} (\pi_ {u} \otimes \pi_ {u}) (c (x)) & = ((1 - \eta_ {u}) \otimes (1 - \eta_ {u})) (c (x)) \\ & = c (x) - (1 \otimes \eta_ {u}) (c (x)) - (\eta_ {u} \otimes 1) (c (x)) + (\eta_ {u} \otimes \eta_ {u}) (c (x)). \end{array}
\]

As \(\varepsilon\) is counit of E,

\[
(1 \otimes \eta_ {u}) (c (x)) = x \otimes u, \quad (\eta_ {u} \otimes 1) (c (x)) = u \otimes x
\]

whence

\[
\left(\eta_ {u} \otimes \eta_ {u}\right) (c (x)) = \left(\eta_ {u} \otimes 1\right) \left(\left(1 \otimes \eta_ {u}\right) (c (x))\right) = \varepsilon (x). u \otimes u;
\]

from this we conclude

\[
\left(\pi_ {u} \otimes \pi_ {u}\right) (c (x)) = c (x) - x \otimes u - u \otimes x + \varepsilon (x). u \otimes u.
\]

On the other hand,

\[
c _ {u} ^ {+} \left(\pi_ {u} (x)\right) = c (x) - x \otimes u - u \otimes x + \varepsilon (x). u \otimes u,
\]

whence formula (4).

As \(\mathrm{E}^{+}\) is a direct factor submodule of \(\mathrm{E}\) , \(\mathrm{E}^{+} \otimes \mathrm{E}^{+}\) can be identified with a direct factor submodule of \(\mathrm{E} \otimes \mathrm{E}\) . With this identification, \(\pi_{u} \otimes \pi_{u}\) is a projector of \(\mathrm{E} \otimes \mathrm{E}\) onto \(\mathrm{E}^{+} \otimes \mathrm{E}^{+}\) . By formula (4), \(c_{u}^{+}\) maps \(\mathrm{E}^{+}\) into \(\mathrm{E}^{+} \otimes \mathrm{E}^{+}\) and \(\pi_{u}\) is a morphism of the cogebra \((\mathrm{E}, c)\) into the cogebra \((\mathrm{E}^{+}, c_{u}^{+})\) .

PROPOSITION 3. If the cogebra \((\mathrm{E}, c)\) is coassociative (resp. cocommutative) (Algebra, Chapter III, § 11, no. 2), so is the cogebra \((\mathrm{E}^{+}, c_{\mathbf{u}}^{+})\) . This follows from the following lemma:

Lemma 1. Let \(\pi: \mathrm{E} \to \mathrm{E}'\) be a surjective cogebra morphism. If \(\mathrm{E}\) is coassociative (resp. cocommutative), so is \(\mathrm{E}'\) .

Let B be an associative K-algebra; the mapping \(f \mapsto f \circ \pi\) is an injective algebra homomorphism of \(\mathrm{Hom}_{\mathbb{K}}(\mathrm{E}', \mathrm{B})\) into \(\mathrm{Hom}_{\mathbb{K}}(\mathrm{E}, \mathrm{B})\) . It then suffices to apply Proposition 1 (resp. Proposition 2) of Algebra, Chapter III, § 11, no. 2.

